
\chapter{Управление отказами}
\label{ch:decline-management}

Когда покупатель видит отказ на странице оплаты, для него всё выглядит
просто: платёж не прошёл. Для платёжной команды на этом простота
заканчивается. Один отказ означает \enquote{подожди и попробуй ещё раз},
другой -- \enquote{эта карта больше не годится}, третий -- \enquote{сломался
маршрут, а не счёт клиента}. Управление отказами и состоит в том, чтобы
различать эти случаи до того, как система начнёт бессмысленно жечь
попытки и раздражать держателя карты.

\section{Анатомия отказа}
\label{sec:decline-anatomy}

На кассе, в приложении и в дашборде все отказы выглядят одинаково:
операция не состоялась. Но рабочая дихотомия \enquote{мягкий / жёсткий отказ} — это мерчантская эвристика. Visa, например, делит decline response codes на категории
\enquote{Issuer will never approve}, \enquote{Issuer cannot approve at this
time}, \enquote{Data quality} и \enquote{Generic}; Mastercard
отдельно запрещает дополнительные CNP-авторизационные запросы для части
decline codes~\cite{visa-core-rules}\cite{mc-tpr}. Поэтому дальше
удобнее говорить не об \enquote{абсолютно истинных} классах, а о том,
какое действие разумно после конкретного ответа.

Мягкий отказ (soft decline) означает, что транзакция отклонена по
временн\'{о}й или устранимой причине --- временный лимит, поведенческий
фильтр, недоступность сервиса. Эмитент как будто говорит:
\enquote{не сейчас, попробуй ещё позже}. К этой группе на практике чаще
всего относят \texttt{51} -- Insufficient Funds, \texttt{61} --
Exceeds Withdrawal
Amount Limit, \texttt{62} -- Restricted Card, \texttt{65} -- Exceeds Withdrawal
Frequency Limit и часть generic-кодов вроде \texttt{05}, если схема вообще
допускает повтор и мерчант остаётся в пределах лимитов
повторов~\cite{visa-core-rules}.

Жёсткий отказ (hard decline) в такой операционной модели означает, что ту
же операцию с теми же данными повторять бессмысленно или прямо запрещено
правилами. Надёжные примеры здесь -- \texttt{04} -- Pick Up, \texttt{14} --
Invalid Account Number, \texttt{41} -- Lost Card и \texttt{43} -- Stolen Card.
Они либо указывают на неверные реквизиты, либо на карту, с которой
продолжать операцию нельзя. Здесь нужен другой платёжный реквизит или ручная проверка~\cite{visa-core-rules}\cite{mc-tpr}.

Есть и коды, которые ломают простое деление. \texttt{54} Expired Card для
Visa относится к категории Data quality: схема прямо трактует
его как случай, где нужно перепроверить платёжную информацию. У
Mastercard правило сформулировано иначе: для CNP мерчант не
должен инициировать дополнительные авторизационные запросы по той же
операции с тем же PAN и сроком действия. Иными словами, проблема здесь в том, что повторять тот же запрос без новых реквизитов нельзя
~\cite{visa-core-rules}\cite{mc-tpr}.

Отдельно стоят технические и маршрутные коды. \texttt{91} -- Issuer
Unavailable и \texttt{96} -- System Malfunction говорят скорее о доступности
эмитента или маршрута, чем о состоянии счёта. Для них схема допускает
повторные попытки в пределах лимитов, но смысл работы здесь другой:
сначала убедиться, что проблема действительно в доступности пути, а не в
самой карте. \texttt{01} -- Refer to Card Issuer тоже не приглашение к
слепому повтору: обычно это знак, что нужен иной канал или
дополнительное действие со стороны держателя либо эмитента
~\cite{visa-core-rules}.

\begin{figure}[tbp]
  \centering
  \includefiguretex[width=\linewidth]{assets/figures/ch25-decline-management/decline-codes}
  \caption{Рабочая классификация кодов ответа по следующему действию.}
  \label{fig:decline-codes}
\end{figure}

Вывод из этой схемы простой: сначала нужно определить обратимость
отказа, и только потом решать, стоит ли повторять попытку. Код ответа
здесь работает не как диагноз, а как сжатая операционная подсказка,
которую всегда нужно сверять с правилами конкретной схемы.

\section{Что на самом деле означают коды ответа}
\label{sec:decline-response-codes}

У платёжной команды всегда есть соблазн прочитать код ответа как
точную причину отказа. Почти всегда это ошибка. Код ответа
(response code) -- это ответ эмитента, но не его объяснение. За одним
и тем же кодом могут стоять десятки причин, и без дополнительных
данных истинную причину обычно не узнать.

\texttt{05} Do Not Honor -- самый распространённый и самый неинформативный
код. Формально он означает \enquote{не проводите операцию}, но почти ничего не
объясняет. За ним могут скрываться антифрод, дневной лимит, несоответствие
поведенческому профилю клиента, временная блокировка в приложении или
внутренняя политика эмитента по конкретному MCC. Часть этих причин
временная и может исчезнуть сама, часть -- постоянная; именно поэтому
\texttt{05} обычно трактуют с осторожностью. У Visa generic decline
response codes сведены в отдельную Category~4 и должны использоваться
только там, где не подходит более точное значение. Уже этого достаточно,
чтобы не читать \texttt{05} как диагноз: это отказ без расшифровки, а не
объяснение причины~\cite{visa-core-rules}.

\texttt{51} Insufficient Funds тоже не сводится к буквальному бытовому
смыслу. Для дебетовых карт он действительно часто означает нехватку
средств на счёте. Для кредитных карт это может быть превышение лимита,
а для некоторых эмитентов -- просто универсальная форма мягкого
отказа, когда они не хотят раскрывать более точную причину.

\texttt{54} Expired Card как раз показывает, почему буквальное чтение опасно.
В Visa Table~7-2 это код класса Data quality, а в Visa Token Service
PAN Lifecycle прямо предусмотрены обновление PAN, срока действия и VAU
updates. Поэтому в recurring- и card-on-file-потоках правильный вопрос
после \texttt{54} часто не \enquote{послать ли ещё раз}, а
\enquote{получили ли мы уже обновлённые реквизиты}. У Mastercard ту же
операционную задачу закрывает Automatic Billing Updater, который
публично позиционируется как способ держать платёжные реквизиты
актуальными
и избегать
decline~\cite{visa-core-rules}\cite{visa-vts}\cite{visa-vau}\cite{mc-bill-pay-solutions}.

\texttt{91} и \texttt{96} -- технические коды, которые не говорят о состоянии
карты или счёта. Они означают, что эмитент или маршрут недоступны, а
значит транзакцию уместно рассматривать как вопрос доступности, а не
кредитного решения. В Visa оба кода попадают в классы, где
повторы разрешены в пределах схемных лимитов; если такие коды
доминируют в статистике, проблема часто лежит в сети,
маршрутизации или доступности хоста. Карта здесь может быть вообще ни при чём.
\cite{visa-core-rules}

\paragraph{Практика.}
Хороший мониторинг отказов сегментирует коды ответа не по двум
группам, а хотя бы по трём-четырём: жёсткие, мягкие финансовые,
мягкие поведенческие и технические. Сводный дашборд с одним числом
доли отказов скрывает структуру: 15\% отказов, половина из которых
приходится на технические коды, означают совсем другую задачу, чем
те же 15\% при преобладании необратимых отказов.

\section{Повторы после отказа: время, частота, каналы}
\label{sec:decline-retry}

Именно на повторах особенно ясно видно, как расходятся интересы
мерчанта и платёжной схемы. Мерчант хочет вернуть конверсию; схема
хочет, чтобы сеть не превращалась в генератор бессмысленных попыток.
Поэтому ошибка в логике повторов обходится потерянной выручкой и мониторингом со штрафами.

Правила схем задают прежде всего внешние границы:
какие коды вообще можно повторять, а какие нельзя. Visa Table~7-2
задаёт категории с разными лимитами повторов для мерчанта; Mastercard для
CNP отдельно запрещает дополнительные авторизационные запросы по той же
транзакции после ряда decline codes, включая \texttt{04},
\texttt{14}, \texttt{41},
\texttt{43} и \texttt{54}~\cite{visa-core-rules}\cite{mc-tpr}.

Отсюда следует важное различие. Схемы почти не говорят мерчанту
\enquote{повтори через 17 минут}; они говорят \enquote{не делай слепой
цикл повторов и не выходи за лимиты}. Если отказ связан с нехваткой
средств или лимитом, серия одинаковых запросов не создаёт деньги и не
расширяет лимит. Если код технический, бессмысленная очередь повторов
лишь засоряет наблюдаемость и повышает риск мониторинга.

Поэтому практическая логика повторов строится из трёх вопросов. Можно
ли повторять этот код по правилам схемы? Нужны ли новые данные -- новая
дата истечения, токен, 3DS, другой платёжный реквизит -- вместо ещё
одной попытки? Если повтор допустим, какой интервал и какой потолок
безопасны для конкретного мерчанта, PSP и вертикали? Mastercard на
публичных продуктовых страницах прямо продаёт это как intelligent retry
strategies~\cite{mc-bill-pay-solutions}.

Канал повтора тоже меняет исход. Visa и Mastercard описывают
токенизацию и оптимизацию одобрений как способы дать эмитенту более
чистый сигнал и убрать лишнюю фрикцию. В этом смысле повтор через
сетевой токен или через более удачный путь маршрутизации действительно
не равен простому повтору PAN-запроса
~\cite{visa-vts}\cite{visa-tokenization-security}\cite{mc-acquirers-psps}\cite{mc-pop}.

\begin{figure}[tbp]
  \centering
  \includefiguretex[width=\linewidth]{assets/figures/ch25-decline-management/decline-retry-tree}
  \caption{Рабочая схема действий после отказа в авторизации.}
  \label{fig:decline-retry-tree}
\end{figure}

Хорошая логика повторов уменьшает не число отказов вообще, а число
бессмысленных повторных попыток. Это разные задачи, и смешивать их
опасно.

\section{Сетевой токен и PAN: дополнительный контекст доверия}
\label{sec:decline-tokens}

Если мерчант работает через сетевой токен, а не хранит PAN,
эмитент получает более управляемый контекст. Visa Token Service
прямо описывает токен как отдельный цифровой идентификатор, с
которым эмитент может задавать token variables, определять авторизованных
token requestors и запрашивать token details, включая device и risk
information. Тот же набор материалов говорит, что токен может быть
ограничен конкретным мобильным устройством, e-commerce-мерчантом или
числом покупок~\cite{visa-vts}\cite{emvco-payment-tokenisation-guide}.

Сетевой токен несёт управляемый контекст о том, где и как реквизит может использоваться. Visa в своих публичных security materials
связывает рост tokenization adoption с ростом approval rates на 2.5\%, а
Mastercard на странице для acquirers и PSP пишет о 3--6~п.п. роста
approval rates у tokenized digital payments
~\cite{visa-tokenization-security}\cite{mc-acquirers-psps}. Эти
цифры принадлежат самим схемам и не обещают магии каждому мерчанту, но
общий вывод устойчив: токен часто даёт эмитенту более чистый сигнал,
чем обычный PAN.

Самый приземлённый выигрыш виден в жизненном цикле реквизитов. У Visa PAN
Lifecycle прямо позволяет обновлять PAN, срок действия и выполнять VAU
updates. У Mastercard Automatic Billing Updater подаётся как способ
держать платёжные реквизиты актуальными и избегать decline.
Поэтому \texttt{54} Expired Card в card-on-file-потоках нередко означает не
\enquote{карта умерла}, а \enquote{у нас устарела ссылка на реквизит}
~\cite{visa-vts}\cite{visa-vau}\cite{mc-bill-pay-solutions}.

На практике миграция на токены откладывается по скучным, но устойчивым причинам.
Нужны интеграция и сертификация; если текущий поток кое-как держится,
миграцию откладывают до заметной боли в виде доли отказов и оттока;
наконец, многие мерчанты вообще не управляют токенами напрямую и
зависят от того, как быстро двигается их PSP. Поэтому разговор о токенах почти всегда охватывает безопасность и операционную зрелость.

Подробнее об архитектуре сетевых токенов~-- в главе \ref{ch:tokenization}.

\section{Маршрутизация через эквайера как инструмент конверсии}
\label{sec:decline-routing}

Маршрутизация через эквайера имеет смысл лишь потому, что контур
приёма не является однотрубной системой. В правилах Mastercard есть
отдельные разделы про authorization routing и domestic transaction
routing, а в публичных материалах Mastercard для acceptance community
routing и payment optimization вынесены в самостоятельный слой
продуктов~\cite{mc-tpr}\cite{mc-acquirers-psps}\cite{mc-pop}.
Уже одно это показывает: путь до эмитента влияет на исход и не
сводится к вопросу \enquote{какой PSP мы подключили}.

Практически из этого следуют две вещи. Во-первых, у разных эквайеров
различаются локальное покрытие, локальные подключения, качество
формирования сообщения и дополнительные сервисы оптимизации.
Во-вторых, появление сетевой оптимизации одобрений показывает,
что иногда правильный ход — более удачный маршрут авторизации или лучше сформированное сообщение
~\cite{mc-pop}.

Но эта свобода не безгранична. Маршрутизация должна оставаться внутри
правил схем, локальной маршрутизации и interchange; она не является
лицензией на маскировку географии мерчанта или искусственное
конструирование тарифного режима. Поэтому легитимная цель routing — повысить долю одобрений и сократить операционные потери~\cite{visa-core-rules}\cite{mc-rules}\cite{mc-tpr}.

В этом смысле smart routing — слой дисциплины поверх платёжного контура: сначала понять ограничения пути,
затем выбрать лучший доступный маршрут, и только потом решать, нужен ли
вообще повтор.

\section{Доля одобрений по вертикалям}
\label{sec:decline-benchmarks}

С долей одобрений всегда одна и та же ловушка: число легко назвать и
трудно понять. 85\% -- это хорошо или тревожно? Ответ зависит от
вертикали, географии, типа карты и самой модели оплаты.

Сначала надо договориться о знаменателе и о самом разрезе. Сетевые правила и мониторинги сами считают approval
performance не одним универсальным числом: Mastercard Transaction
Processing Rules, например, упоминают Minimum Average Approval Rate для
конкретных cross-border- и card-not-present-контекстов. Уже это
предупреждает против сравнения всех потоков одной агрегированной
метрикой~\cite{mc-tpr}.

Но сами диапазоны стареют быстрее, чем успевает выйти книга. Они
зависят от вертикали, страны, доли 3DS, качества токенизации, текущих
волн мошенничества и даже календарного окна. Поэтому для печатной
главы важнее понять, как читать эти цифры, чем пытаться запомнить
конкретные проценты.

Сравнивать розницу с сервисами путешествий, подписки с разовыми
интернет-платежами, а domestic-поток с cross-border в одной колонке
обычно бессмысленно. В recurring отдельно добавляется контекст
сохранённого реквизита: если цепочка согласия не оформлена по
правилам SCF, эмитент
видит обычный CNP-запрос, а не продолжение уже данного согласия
\cite{visa-scf}\cite{mc-scf}.

\paragraph{Важно.}
Долю одобрений нельзя читать одним числом по продукту. Метод только
один: смотреть разрезы по странам, вертикалям и типам интеграции,
иначе реальные причины отказов прячутся внутри агрегата.

\paragraph{Практика.}
Измеряйте долю одобрений раздельно по сегментам: новые и вернувшиеся
клиенты, карты с 3DS и без неё, PAN и токен, внутренний рынок и
трансграничные операции. Агрегированная метрика почти бесполезна для
оптимизации. Сегментация быстро показывает, где сосредоточены
отказы: как правило, основную массу проблем даёт не весь поток, а
сравнительно узкий класс транзакций с особыми характеристиками.

\section*{Итоги главы}
\begin{itemize}
  \item Разделение на мягкие и жёсткие отказы полезно как операционная
    модель, но правила схем устроены тоньше: у Visa есть
    категории never approve / cannot approve at this time / data
    quality / generic, а Mastercard отдельно запрещает часть
    повторов той же CNP-операции.
  \item \texttt{05} Do Not Honor не объясняет причину отказа, а \texttt{54}
    Expired Card не всегда означает \enquote{карта навсегда
    недействительна}: в card-on-file-потоке это часто вопрос
    обновлённых реквизитов, токенов и updater-сервисов.
  \item Visa Transaction Retry Rules and Excessive Attempt Restrictions
    и Mastercard Excessive Attempts Rules задают внешние границы
    повторов; схемы ограничивают число попыток и типы отказов, для которых повтор допустим.
  \item Сетевой токен (network token) даёт эмитенту более управляемый
    контекст, а у Visa и Mastercard публичные
    материалы прямо связывают tokenization с ростом approval rate;
    это улучшение сигнала.
  \item Маршрутизация и payment optimization — реальный слой acceptance stack; использовать его нужно внутри правил схем, маршрутизации и interchange.
  \item Долю одобрений имеет смысл читать только в разрезах:
    география, тип транзакции, контекст сохранённого реквизита, PAN или
    токен, внутренний рынок или трансграничный поток; одно
    агрегированное число почти всегда скрывает разные механизмы
    отказов.
\end{itemize}
